\documentclass[11pt,a4paper]{article}
\usepackage[T1]{fontenc}
\usepackage{isabelle,isabellesym}

\usepackage{amsmath}
\usepackage{amssymb}
\usepackage{stmaryrd}

% this should be the last package used
\usepackage{pdfsetup}

\urlstyle{rm}
\isabellestyle{it}

\begin{document}

\title{Comparison of Viscosity Solutions for Geometric Operators}
\author{Dmitriy Traytel}
\maketitle

\begin{abstract}
  We prove comparison principles for viscosity solutions of the equation
  $F(Du, D^2u) = 1$, where $F$ is any second-order operator on
  $\mathbb{R}^n \times \mathbb{R}^{n \times n}$ with ten properties,
  collected in a locale of geometric operators: degenerate ellipticity,
  positive homogeneity in the Hessian, invariance under positive dilations
  of the gradient and under rotations, a bound on the effect of shifting
  the Hessian by a multiple of the identity, agreement with both
  semicontinuous envelopes where the gradient does not vanish and with the
  lower one where it does, and the upper envelope $F^*$ below $1$ at
  $(0, 0)$ and uniformly at $(q, -\delta I)$.  The operator may thus be
  discontinuous where the gradient vanishes, as operators of geometric
  equations are.  The entry defines the semicontinuous envelopes of an
  operator and their invariances, and proves the maximum principle for an
  upper semicontinuous subsolution and a lower semicontinuous
  supersolution on a bounded set, by doubling the variables in the
  sup-convolutions of the two functions, with a penalty whose gradient
  stays away from zero; comparison between two domains; and comparison
  and uniqueness on expandable sets and on compact sets.  The viscosity
  notions are those of \texttt{Second\_Order\_Viscosity\_Analysis}, at
  local touching and at touching over a set, and the second-order tools
  (the sup-convolution, Jensen's lemma and the doubling of variables) come
  from there as well.  We follow the User's Guide of
  Crandall, Ishii and Lions~\cite{CrandallIshiiLions}.
\end{abstract}

\tableofcontents

\parindent 0pt\parskip 0.5ex

\input{session}

\bibliographystyle{abbrv}
\bibliography{root}

\end{document}
